\subsection{CONNECTED SPACES AND CONNECTED SETS}

DEFINITION 1. A topological space X is said to be connected if it is not the union of two disjoint non-empty open sets.

An equivalent definition is obtained by replacing the words ``open sets'' by ``closed sets''. X is connected if and only if the only subsets of X which are both open and closed are the empty set and the whole space X.

If X is connected and if A, B are two non-empty open (resp. closed) subsets such that \(A \cup B = X\) , then \(A \cap B \neq \varnothing\) .

Examples. * 1) We shall see in Chapter IV, § 2, no. 5 that the real line is connected, and that the rational line is not. *

\begin{enumerate}
\def\labelenumi{\arabic{enumi})}
\setcounter{enumi}{1}
\tightlist
\item
  A discrete space which has more than one point is not connected.
\end{enumerate}

Observe that if \((\mathbf{U}_{t})_{t\in\mathbf{I}}\) is a partition of a topological space X consisting of open sets {[}necessarily non-empty, from the definition of a partition (Set Theory, R, § 4, no. 4){]} then each of the \(U_{t}\) is both open and closed in X, for the complement of \(U_{t}\) is the union of the \(U_{x}\) for which \(x\neq t\) . The open subsets of X are the subsets A such that \(A\cap U_{t}\) is open in \(U_{t}\) for each \(t\in I\) , so that X can be identified with the sum of the \(U_{t}\) (§ 2, no. 4, Example 3) and is not connected if I has more than one element.

DEFINITION 2. A subset A of a topological space X is said to be a connected set if the subspace A of X is connected.

For A to be a connected subset of X it is necessary and sufficient that, for each covering of A by two open (or closed) subsets B, C of X such that \(A \cap B\) and \(A \cap C\) are non-empty, we have \(A \cap B \cap C \neq \emptyset\) .

Examples. In any topological space, the empty set and every set consisting of a single point are connected. In a Hausdorff space X, every finite set consisting of more than one point is not connected, and more generally every subset of X which has more than one point and which has at least one isolated point is not connected.

If a dense subset A is connected, then the whole space X is connected; otherwise there would exist two non-empty disjoint open subsets M, N of X such that \(M \cup N = X\) , and \(M \cap A\) , \(N \cap A\) would be two disjoint non-empty open subsets of A whose union is A. Hence we have

PROPOSITION 1. If A is a connected set, then every set B such that \(A \subset B \subset \overline{A}\) is connected.

PROPOSITION 2. The union of a family of connected sets whose intersection is non-empty is connected.

Let \((\mathbf{A}_t)_{i\in I}\) be a family of connected subsets of \(\mathbf{X}\) , all of which contain the same point \(x\) ; we have to show that

\[
\mathbf {A} = \bigcup_ {i} \mathbf {A} _ {i}
\]

is connected. If not, there are two open sets B and C such that \(B \cap A\) and \(C \cap A\) are non-empty, and \(A \subset B \cup C\) and \(A \cap B \cap C = \varnothing\) . x belongs to one of the sets B, C, say \(x \in B\) ; on the other hand one of the sets \(A_t\) , say \(A_x\) , meets C; we have therefore \(A_x \subset B \cup C\) , \(A_x \cap B \cap C = \varnothing\) and \(B \cap A_x\) and \(C \cap A_x\) are non-empty. Hence \(A_x\) is not connected, which is a contradiction.

COROLLARY. Let \((\mathbf{A}_n)_{n\geqslant 0}\) be an infinite sequence of connected sets such that \(\mathbf{A}_{n + 1}\cap \mathbf{A}_n\neq \emptyset\) for all \(n\geqslant 0\) . Then the union \(\bigcup_{n = 0}^{\infty}\mathbf{A}_n\) is connected.

By induction on \(n\) we see immediately that the set \(\mathbf{B}_n = \bigcup_{i=0}^n \mathbf{A}_i\) is connected for all \(n\) , by Proposition 2. The sets \(\mathbf{B}_n\) have a non-empty intersection; hence their union, equal to \(\bigcup_{n=0}^{\infty} \mathbf{A}_n\) , is connected by Proposition 2.

PROPOSITION 3. Let A be a subset of a topological space X. If B is a connected subset of X which meets both A and \(\mathbb{C}\) A, then B meets the frontier of A.

For otherwise the intersections of B with the interior and exterior of A would be two open subsets of B which form a partition of B, and B would not be connected.

COROLLARY. In a connected space X, every non-empty set other than X itself has at least one frontier point.

\subsection{IMAGE OF A CONNECTED SET UNDER A CONTINUOUS MAPPING}

PROPOSITION 4. Let A be a connected subset of a topological space X, and let f be a continuous mapping of X into a topological space X'. Then \(f(\mathbf{A})\) is connected.

Suppose \(f(\mathbf{A})\) is not connected. Then there exist two sets \(\mathbf{M}'\) , \(\mathbf{N}'\) which are open in \(f(\mathbf{A})\) and which form a partition of \(f(\mathbf{A})\) ; hence \(\mathbf{A} \cap \overline{f}^{1}(\mathbf{M}')\) and \(\mathbf{A} \cap \overline{f}^{1}(\mathbf{N}')\) are open in \(\mathbf{A}\) and form a partition of \(\mathbf{A}\) ; this contradicts the hypothesis that \(\mathbf{A}\) is connected.

The inverse image of a connected set under a continuous mapping need not be connected; consider for example a mapping of a discrete space into a space consisting of one point.

From Proposition 4 we derive another characterization of non-connected spaces:

PROPOSITION 5. For a topological space X to be not connected it is necessary and sufficient that there exists a surjective continuous mapping of X onto a discrete space containing more than one point.

The condition is sufficient by Proposition 4. Conversely, if X is not connected, there exist two non-empty disjoint open subsets A, B whose union is X, and the mapping f of X onto a discrete space of two elements \(\{a, b\}\) , defined by \(f(A) = \{a\}\) and \(f(B) = \{b\}\) , is continuous.

\subsection{QUOTIENT SPACES OF A CONNECTED SPACE}

PROPOSITION 6. Every quotient space of a connected space is connected.

This is an immediate consequence of Proposition 4 of no. 2.

PROPOSITION 7. Let X be a topological space and R an equivalence relation on X. If the quotient space X/R is connected, and if each equivalence class mod R is connected, then X is connected.

Suppose X is not connected. Then there is a partition of X into two open sets A, B. The sets A, B are saturated with respect to R; for if \(x \in A\) then the equivalence class M of x cannot meet B, otherwise the sets \(A \cap M\) , \(B \cap M\) would form a partition of M into two sets open in M, which is impossible since M is connected. The canonical images of A and B are therefore open sets in X/R and form a partition of X/R; this contradicts the assumption that X/R is connected.

\subsection{PRODUCT OF CONNECTED SPACES}

PROPOSITION 8. Every product of connected spaces is connected. Conversely, if a product of non-empty spaces is connected, then each of the factors is connected.

Let \(\mathbf{X} = \prod_{i \in I} \mathbf{X}_i\) be a product of topological spaces. If the \(\mathbf{X}_i\) are non-empty, we have \(\mathbf{X}_i = \text{pr}_i\mathbf{X}\) for each \(i \in I\) ; hence if \(\mathbf{X}\) is connected so are the \(\mathbf{X}_i\) (no. 2, Proposition 4). Conversely, suppose that each of the \(\mathbf{X}_i\) is connected and \(\mathbf{X}\) is not. By Proposition 5 of no. 2, there exists a continuous surjective mapping \(f: \mathbf{X} \to \mathbf{X}'\) , where \(\mathbf{X}'\) is a discrete space which contains more than one point. Let \(a = (a_i)\) be any point of \(\mathbf{X}\) , and \(x\) any index; the partial mapping \(f_x: \mathbf{X}_x \to \mathbf{X}'\) , defined by \(f_x(x) = f((y_i))\) where \(y_x = x\) and \(y_i = a_i\) if \(i \neq x\) , is continuous on \(\mathbf{X}_x\) ; since \(\mathbf{X}_x\) is connected, \(f_x\) must be constant on \(\mathbf{X}_x\) . It follows immediately by induction that \(f(x) = f(a)\) for all points \(x = (x_i)\) such that \(x_i = a_i\) for all but a finite number of indices \(i \in I\) . But these points \(x\) form a dense subset of \(\mathbf{X}\) ( \(\S 4\) , no. 3, Proposition 8). Hence \(f\) is continuous on \(\mathbf{X}\) and constant on a dense subset of \(\mathbf{X}\) , and therefore constant on \(\mathbf{X}\) ( \(\S 8\) , no. 1, Proposition 2, Corollary 1). But this contradicts the definition of \(f\) .

\subsection{COMPONENTS}

Given a point x of a topological space X, the union of the connected subsets of X which contain x is connected (no. 1, Proposition 2); it is therefore the largest connected subset of X which contains x.

DEFINITION 3. The component (or connected component) of a point of a topological space X is the largest connected subset of X which contains this point. The components of a subset A of X are the components of the points of A, relative to the subspace A of X.

If a space is connected, the component of each point is the whole space. If a space X is such that for each pair \((x, y)\) of points of X there is a connected set containing x and y, then X is connected.

A space X is said to be totally disconnected if the component of each point of X consists of the point alone. A subset A of X is a totally disconnected set if the subspace A of X is totally disconnected.

A discrete space is totally disconnected, but one should beware of confusing these two notions; for example, we shall see in Chapter IV, § 2, no. 5, that the rational line, which is not a discrete space, is totally disconnected.

\pandocbounded{\includegraphics[keepaspectratio]{images/part001_10033d49.jpg}}

A set which is both open and closed contains the component of each of its points, whence the component of a point is contained in the intersection of the sets which are both open and closed and which contain this point. However, the component of a point is not necessarily equal to this intersection (cf.~Exercise 9 and Chapter II, § 4, no. 4, Proposition 6).

PROPOSITION 9. The component of any point in a topological space X is a closed set. The relation ``y belongs to the component of x'' is an equivalence relation R\{x, y\} on X, and the equivalence classes are the components of X. The quotient space X/R is totally disconnected.

The first part of the proposition is an immediate consequence of Definition 3 and the fact that the closure of a connected set is connected (no. 1, Proposition 1). Since the union of connected sets which have a point in common is connected (no. 1, Proposition 2), the relation R is transitive, hence is an equivalence relation (since it is obviously reflexive and symmetric) and the equivalence class of x with respect to R is the component of x. It remains to show that X/R is totally disconnected. Let \(f: X \to X/R\) be the canonical mapping, and let F be a closed set in X/R containing at least two distinct points; the inverse image \(\overline{f}^{1}(F)\) of F is closed in X, saturated with respect to R, and contains at least two distinct components of X and hence is not connected. Hence there exist two non-empty closed sets B, C in X such that \(B \cap C = \varnothing\) and \(B \cup C = \overline{f}^{1}(F)\) . The component of any point x of \(\overline{f}^{1}(F)\) in \(\overline{f}^{1}(F)\) is the same as the component of x in X (by the definition of R) and therefore B and C, which are both open and closed in \(\overline{f}^{1}(F)\) , are saturated with respect to R. Hence \(f(B)\) and \(f(C)\) are closed in X/R, and \(f(B) \cup f(C) = F\) and \(f(B) \cap f(C) = \varnothing\) ; this shows that F is not connected and consequently that X/R is totally disconnected.

PROPOSITION 10. In a product space \(\mathbf{X} = \prod_{i\in J}\mathbf{X}_i\) , the component of \(x = (x_i)\) in \(\mathbf{X}\) is the product of the components of \(x_i\) in the factors \(\mathbf{X}_i\) .

This product set is connected (no. 4, Proposition 8). Conversely, if A is a connected subset of X which contains x, then \(\mathrm{pr}_{i}(A)\) is a connected set (no. 2, Proposition 4) which contains x; since \(A \subset \prod_{i} \mathrm{pr}_{i}(A)\) , it follows that A is contained in the product of the components of the \(x_{i}\) .

\subsection{LOCALLY CONNECTED SPACES}

DEFINITION 4. A topological space X is said to be locally connected if each point of X has a fundamental system of connected neighbourhoods.

* We shall see in Chapter IV, § 2, no. 5, that the real line is a locally connected space. *

The existence, at each point x of a space X, of one connected neighbourhood of x by no means implies that X is locally connected. In particular, X can be connected but not locally connected (Exercises 2 and 13). Conversely, a space can be locally connected but not connected (e.g.~a discrete space which contains more than one point).

PROPOSITION 11. A necessary and sufficient condition for a space X to be locally connected is that every component of an open set in X is open in X.

The condition is sufficient, since the component of x relative to an open neighbourhood of x is then a neighbourhood of x in X.

Conversely, let A be an open subset of a locally connected space X, let B be a component of A, and let \(x \in B\) . Let V be a connected neighbourhood of x contained in A; by the definition of components, V is contained in B; hence B is open in X (§ 1, no. 2, Proposition 1).

The components of a locally connected space X therefore form a partition of X into open sets, and hence X is the sum (§ 2, no. 4, Example 3) of its components.

COROLLARY. Let U be an open subset of a locally connected space X, and let V be a component of U. Then the frontier of V (relative to X) is contained in the frontier of U.

For V is open and closed in U, hence a frontier point of V (relative to X) cannot belong to U, for it would also be a frontier point of V relative to U, and there is none.

PROPOSITION 12. Every quotient space of a locally connected space is locally connected.

Let X be a locally connected space, R an equivalence relation on X, \(\varphi: X \to X/R\) the canonical mapping. Let A be an open subset of

X/R and C a component of A. Then \(\overline{\varphi}^{1}(C)\) is a union of components of \(\overline{\varphi}^{1}(A)\) ; for if \(x \in \overline{\varphi}^{1}(C)\) and if K is the component of \(x\) in \(\overline{\varphi}^{1}(A)\) then \(\varphi(K)\) is connected (no. 2, Proposition 4), is contained in A, and contains \(\varphi(x)\) ; hence \(\varphi(K) \subset C\) by the definition of C, and therefore \(K \subset \overline{\varphi}^{1}(C)\) . Since X is locally connected and \(\overline{\varphi}^{1}(A)\) is open in X, it follows from Proposition 11 that \(\overline{\varphi}^{1}(C)\) is open in X; consequently C is open in X/R and hence, by Proposition 11 again, X/R is locally connected.

PROPOSITION 13. a) Let \((\mathbf{X}_t)_{t \in I}\) be a family of locally connected spaces such that \(\mathbf{X}_t\) is connected for all but a finite number of indices \(t \in I\) . Then the product space \(\mathbf{X} = \prod_{i \in I} \mathbf{X}_t\) is locally connected.

\begin{enumerate}
\def\labelenumi{\alph{enumi})}
\setcounter{enumi}{1}
\item
  Conversely, if the product of a family \((\mathbf{X}_{t})\) of non-empty topological spaces is locally connected, then each \(X_{t}\) is locally connected, and \(X_{t}\) is connected for all but a finite number of indices.
\item
  Let \(\mathbf{J}\) be the finite subset of \(\mathbf{I}\) such that \(\mathbf{X}_i\) is not connected if and only if \(i \in \mathbf{J}\) . Let
\end{enumerate}

\[
\mathrm{U} = \prod_ {i \in I} \mathrm{U} _ {i}
\]

be an elementary set containing a point \(x = (x_{t})\) of \(\mathbf{X}\) and let \(\mathbf{K}\) be the finite subset of \(\mathbf{I}\) such that \(\mathbf{U}_{t} \neq \mathbf{X}_{t}\) if and only if \(t \in \mathbf{K}\) . Let \(\mathbf{V}_{t}\) be \(\mathbf{X}_{t}\) for \(t \notin J \cup K\) , and let \(V_{t}\) be a connected neighbourhood of \(x_{t}\) contained in \(U_{t}\) for \(t \in J \cup K\) ; then

\[
\mathbf {V} = \prod_ {i \in I} \mathbf {V} _ {i}
\]

is connected (by Proposition 8 of no. 4) and is a neighbourhood of x contained in U. Hence X is locally connected.

\begin{enumerate}
\def\labelenumi{\alph{enumi})}
\setcounter{enumi}{1}
\tightlist
\item
  Let \(a = (a_{i})\) be a point of \(\mathbf{X}\) and let \(\mathbf{V}\) be a connected neighbourhood of \(a\) in \(\mathbf{X}\) . Since we have \(\mathrm{pr}_{\mathfrak{i}}\mathbf{V} = \mathbf{X}_{\mathfrak{i}}\) except for a finite number of indices (§ 4, no. 1) it follows from no. 2, Proposition 4 that the \(\mathbf{X}_{\mathfrak{i}}\) are connected, for all but a finite number of indices. On the other hand, for each \(x \in I\) , each \(a_{x} \in X_{x}\) and each neighbourhood \(V_{x}\) of \(a_{x}\) in \(X_{x}\) , there is a point \(x\) of \(\mathbf{X}\) such that \(\mathrm{pr}_{\mathfrak{x}}x = a_{x}\) , and
\end{enumerate}

\[
\mathbf {V} = \mathbf {V} _ {x} \times \prod_ {i \neq x} \mathbf {X} _ {i}
\]

is a neighbourhood of \(x\) in \(\mathbf{X}\) ; \(\mathbf{V}\) therefore contains a connected neighbourhood \(W\) of \(x\) , whose projection \(\mathrm{pr}_xW\) is a connected neighbourhood of \(a_x\) contained in \(V_x\) (no. 2, Proposition 4 and § 4, no. 2, Proposition 5). Hence each \(X_x\) is locally connected.

\subsection{APPLICATION : THE POINCARÉ-VOLTERRA THEOREM}

THEOREM 1. Let X be a topological space satisfying axiom (O \(_{III}\) ) (but not necessarily Hausdorff), and suppose X is connected and locally connected. Let Y be a topological space whose topology has a countable base, and let p: X → Y be a continuous mapping such that, for each y ∈ Y, \(\overline{p}^{1}(y)\) is a discrete subspace of X. Finally let B be a set of subsets of X whose interiors cover X and such that:

\begin{enumerate}
\def\labelenumi{(\roman{enumi})}
\tightlist
\item
  The restriction of \(p\) to each \(V \in \mathfrak{B}\) is a closed mapping of \(V\) into \(Y\) .\\
\item
  Every \(V \in \mathfrak{B}\) has a countable subset which is dense in \(V\) .
\end{enumerate}

Then the space X is the union of a countable family of open sets each of which is contained in a set of B.

Let \(\mathfrak{B}\) be a countable base of the topology of Y. We shall say that a pair (W, U) is distinguished if (i) \(U \in \mathfrak{B}\) and (ii) W is a component of \(\overline{p}(U)\) contained in a set of \(\mathfrak{B}\) .

LEMMA 1. If \(x\) is any point of \(\mathbf{X}\) , there is a distinguished pair \((\mathbf{W}, \mathbf{U})\) such that \(x \in \mathbf{W}\) .

The inverse image \(\overline{p} (p(x))\) is discrete and therefore there is a neighbourhood of \(x\) in \(\mathbf{X}\) all of whose points \(x^{\prime}\) other than \(x\) have an image \(p(x^{\prime})\neq p(x)\) ; since \(\mathbf{X}\) satisfies \((\mathrm{O}_{\mathrm{III}})\) , there is a closed neighbourhood \(\mathbf{V}\) of \(x\) with this property, and we may assume also that \(\mathbf{V}\) is contained in a set of \(\mathfrak{B}\) . Let \(\mathbf{F}\) be the frontier of \(\mathbf{V}\) in \(\mathbf{X}\) . By condition (i) of the theorem, \(p(\mathbf{F})\) is closed in \(\mathbf{Y}\) ; and since \(p(\mathbf{F})\) does not contain \(p(x)\) , there is a set \(\mathbf{U}\in \mathfrak{B}\) which contains \(p(x)\) and does not meet \(p(\mathbf{F})\) . Let \(\mathbf{W}\) be the component of \(x\) in \(\overline{p}^{1}(\mathbf{U})\) ; then it is enough to show that \(\mathbf{W}\subset \mathring{\mathbf{V}}\) . If this were not so, then \(\mathbf{W}\) would meet \(\mathbf{F}\) (no. 1, Proposition 3) and therefore \(p(\mathbf{F})\) would meet \(\mathbf{U}\) , contrary to the definition of \(\mathbf{U}\) .

LEMMA 2. If (W, U) is a distinguished pair then the set of all distinguished pairs (W', U') such that W' meets W is countable.

Since \(\mathfrak{B}\) is countable it is enough to show that, given \(\mathbf{U}' \in \mathfrak{B}\) , the set of distinguished pairs \((\mathbf{W}', \mathbf{U}')\) such that \(\mathbf{W}'\) meets \(\mathbf{W}\) is countable. Now these sets \(\mathbf{W}'\) are open, since \(\mathbf{X}\) is locally connected (no. 6, Proposition 11) and mutually disjoint since they are components of \(\overline{\boldsymbol{p}}(\mathbf{U}')\) ; hence the sets \(\mathbf{W}' \cap \mathbf{W}\) are open and mutually disjoint. But \(\mathbf{W}\) contains a countable subset which is dense in \(\mathbf{W}\) ; hence the set of \(\mathbf{W}'\) such that \(\mathbf{W}' \cap \mathbf{W}\) is not empty is also countable.

To prove Theorem 1, consider the following relation R between two points \(x\) , \(x'\) of X: ``There exists a finite sequence of distinguished pairs \((W_i, U_i)\) ( \(1 \leqslant i \leqslant n\) ) such that \(x \in W_1\) and \(x' \in W_n\) and that \(W_i \cap W_{i+1} \neq \emptyset\) for \(1 \leqslant i \leqslant n - 1\) .''

Lemma 1 states that R is reflexive, and it is readily verified that R is symmetric and transitive, so that R is an equivalence relation; also, since the \(W_{i}\) are open, every equivalence class mod R is open in X. But X is connected; hence there can be only one equivalence class, i.e.~any two points of X are congruent mod R. We shall deduce from this that X is the union of a countable family of first elements of distinguished pairs, and this will prove Theorem 1. For this, let x be any point of X, and define by induction on n a sequence \((\mathbf{C}_{n})\) of open subsets of X as follows: by Lemma 1 there is a distinguished pair \((\mathbf{W}_{1}, \mathbf{U}_{1})\) such that \(x \in W_{1}\) , and we take \(C_{1} = W_{1}\) ; if n \textgreater{} 1 then \(C_{n}\) is to be the union of all first elements W of distinguished pairs (W, U) such that W meets \(C_{n-1}\) . By induction on n one shows immediately, by virtue of Lemma 2, that \(C_{n}\) is a countable union of first elements of distinguished pairs. Finally, every \(x' \in X\) belongs to some \(C_{n}\) ; for there is a finite sequence

\[
\left(\mathrm{W} _ {i} ^ {\prime}, \mathrm{U} _ {i} ^ {\prime}\right) _ {1 \leqslant i \leqslant m}
\]

of distinguished pairs such that \(x \in W_1'\) , \(x' \in W_m'\) and \(W_i' \cap W_{i+1}' \neq \emptyset\) for \(1 \leqslant i \leqslant m - 1\) , and by induction on \(i\) we see that \(W_i' \subset C_{i+1}\) for all \(i\) , so that \(x' \in C_{m+1}\) .

Q.E.D.

COROLLARY 1. Let Y be a regular space whose topology has a countable base (*). Let X be a connected and locally connected space, and let \(p: \mathbf{X} \to \mathbf{Y}\) be a continuous mapping with the following property: for each \(x \in \mathbf{X}\) there is a closed neighbourhood V of \(x\) in X such that the restriction of \(p\) to V is a homeomorphism of V onto a closed subspace of Y. Then X is regular and the topology of X has a countable base.

First, the hypotheses imply that X is regular (§ 8, no. 4, Proposition 13). Let us show that the conditions of Theorem 1 are satisfied if we take B to be the set of all closed subsets V of X such that the restriction of p to V is a homeomorphism of V onto a closed subspace of Y. By hypothesis, the interiors of the sets of B cover Y and, by virtue of the assumption on Y, each \(V \in B\) has a countable base and therefore contains a countable dense subset (§ 1, no. 6, Proposition 6). Furthermore, if \(x \in \overline{p}(y)\) , there is a neighbourhood \(V \in B\) of x in X such that V contains no point of \(\overline{p}(y)\) other than \(x\) , and hence \(\overline{p}(y)\) is a discrete space. We may therefore apply Theorem 1, which shows that X is the union of a countable family \((\mathrm{T}_n)_{n\geqslant 0}\) of open sets, such that each subspace \(\mathrm{T}_n\) has a countable base \((\mathrm{U}_{mn})_{m\geqslant 0}\) . Then the \(\mathrm{U}_{mn}\) ( \(m\geqslant 0\) , \(n\geqslant 0\) ) form a base of the topology of X (§ 3, no. 1, Remark).

COROLLARY 2. Let X be locally compact, connected and locally connected, and suppose each point of X has a neighbourhood which has a countable base. Let Y be a Hausdorff space whose topology has a countable base, and let \(p: \mathbf{X} \to \mathbf{Y}\) be a continuous mapping such that, for each \(y \in \mathbf{Y}\) , \(\overline{p}^1(y)\) is a discrete subspace of X. Then the topology of X has a countable base.

For each \(x \in X\) , let \(V_{x}\) be a compact neighbourhood of x in X which has a countable base. It follows from §9, no. 4, Theorem 2, Corollary 2, that the set B of the \(V_{x}\) satisfies the conditions of Theorem 1, and we complete the proof as in Corollary 1.

Notice that, in this corollary, it can happen that the restriction of p to an arbitrarily small neighbourhood V of a point of X is not a homeomorphism of V onto \(p(\mathrm{V})\) .

COROLLARY 3 (Poincaré-Volterra Theorem). Let Y be a locally compact, locally connected space whose topology has a countable base. Let X be a connected Hausdorff space, and let \(p: \mathbf{X} \to \mathbf{Y}\) be a continuous mapping which has the following property: for each \(x \in \mathbf{X}\) there is an open neighbourhood U of \(x\) in X such that the restriction of \(p\) to U is a homeomorphism of U onto an open subspace of Y. Then X is locally compact and locally connected, and the topology of X has a countable base.

Clearly X is locally connected. Also each \(x \in X\) has an open neighbourhood U in X such that the restriction of p to U maps U homeomorphically onto an open subspace \(p(U)\) of Y. Since \(p(U)\) is a locally compact subspace of Y ( \(§ 9, no. 7, Proposition 13\) ), there is a compact neighbourhood W of \(p(x)\) contained in \(p(U)\) , whence \(U \cap \overline{p}(W)\) is a compact neighbourhood of x contained in U; thus X is locally compact, since by hypothesis X is Hausdorff. \(U \cap \overline{p}(W)\) , being compact, is closed in X ( \(§ 9, no. 3, Proposition 4\) ) and the conditions of Corollary 1 are therefore satisfied; hence the topology of X has a countable base.

\BourbakiExercises
\BourbakiExerciseGroup

\begin{enumerate}
\def\labelenumi{\arabic{enumi})}
\item
  Find all possible topologies on a set of two elements or three elements.
\item
  \begin{enumerate}
  \def\labelenumii{\alph{enumii})}
  \tightlist
  \item
    Let X be an ordered set. Show that the set of intervals
  \end{enumerate}
\end{enumerate}

\[
[ x, \rightarrow [ \quad (\text { resp. } ] \leftarrow , x ])
\]

is a base of a topology on X; this topology is called the right (resp. left) topology of X. In the right topology, any intersection of open sets is an open set, and the closure of \(\{x\}\) is the interval {]}\(\leftarrow\) , x{]}.

\begin{enumerate}
\def\labelenumi{\alph{enumi})}
\setcounter{enumi}{1}
\item
  A topological space is said to be a Kolmogoroff space if it satisfies the following condition: given any two distinct points \(x\) , \(x'\) of \(X\) , there is a neighbourhood of one of these points which does not contain the other. Show that an ordered set with the right topology is a Kolmogoroff space.
\item
  Let X be a Kolmogoroff space in which every intersection of open sets is an open set. Show that \(x \in \overline{\{x'\}}\) is an order relation between x and \(x'\) in X, and that if this relation is written as \(x \leqslant x'\) then the given topology on X is identical with the right topology determined by this ordering.
\item
  Deduce from c) that if X is a Kolmogoroff space then every finite non-empty subset of X has at least one isolated point. Hence, if X has no isolated point, every non-empty open subset of X is infinite.
\end{enumerate}

\begin{enumerate}
\def\labelenumi{\arabic{enumi})}
\setcounter{enumi}{2}
\tightlist
\item
  If \(\mathbf{A}\) is any subset of a topological space \(\mathbf{X}\) , let \(\alpha(\mathbf{A})\) denote \(\overline{\mathbf{A}}\) and \(\beta(\mathbf{A})\) denote \(\overline{\mathbf{A}}\) . Then \(\mathbf{A} \subset \mathbf{B}\) implies \(\alpha(\mathbf{A}) \subset \alpha(\mathbf{B})\) and \(\beta(\mathbf{A}) \subset \beta(\mathbf{B})\) . \(a)\) Show that if \(\mathbf{A}\) is open then \(\mathbf{A} \subset \alpha(\mathbf{A})\) , and that if \(\mathbf{A}\) is closed then \(\beta(\mathbf{A}) \subset \mathbf{A}\) .
\end{enumerate}

\begin{enumerate}
\def\labelenumi{\alph{enumi})}
\setcounter{enumi}{1}
\item
  Deduce from a) that if A is any subset of X, then \(\alpha(\alpha(A)) = \alpha(A)\) and \(\beta(\beta(A)) = \beta(A)\) .
\item
  Give an example of a set A \(*\) (on the real line) \(_{*}\) such that the seven sets A, \(\mathring{A}\) , \(\overline{A}\) , \(\alpha(A)\) , \(\beta(A)\) , \(\beta(\overline{A})\) , \(\alpha(\mathring{A})\) are all distinct and satisfy no relations of inclusion except the following: \(\mathring{A} \subset A \subset \overline{A}\) , \(\mathring{A} \subset \alpha(\mathring{A}) \subset \beta(A) \subset \overline{A}\) , \(\mathring{A} \subset \alpha(A) \subset \beta(\overline{A}) \subset \overline{A}\) , \(\alpha(\mathring{A}) \subset \alpha(A)\) , \(\beta(A) \subset \beta(\overline{A})\) .
\item
  Show that if U, V are two open sets such that \(U \cap V = \varnothing\) , then also \(\alpha(U) \cap \alpha(V) = \varnothing\) {[}use b){]}.
\end{enumerate}

\begin{enumerate}
\def\labelenumi{\arabic{enumi})}
\setcounter{enumi}{3}
\tightlist
\item
  \begin{enumerate}
  \def\labelenumii{\alph{enumii})}
  \tightlist
  \item
    Give an example * on the real line * of two open sets A, B such that the four sets A ∩ B, B ∩ A̅, A ∩ B and A̅ ∩ B are all distinct.
  \end{enumerate}
\end{enumerate}

\begin{enumerate}
\def\labelenumi{\alph{enumi})}
\setcounter{enumi}{1}
\tightlist
\item
  * Give an example of two intervals A, B on the real line, such that A ∩ B is not contained in A ∩ B. *
\end{enumerate}

\begin{enumerate}
\def\labelenumi{\arabic{enumi})}
\setcounter{enumi}{4}
\tightlist
\item
  If A is any subset of a topological space X, the frontier of A is denoted by \(\operatorname{Fr}(\mathbf{A})\) .
\end{enumerate}

\begin{enumerate}
\def\labelenumi{\alph{enumi})}
\item
  Show that \(\operatorname{Fr}(\overline{\mathbf{A}}) \subset \operatorname{Fr}(\mathbf{A})\) , \(\operatorname{Fr}(\mathring{\mathbf{A}}) \subset \operatorname{Fr}(\mathbf{A})\) , and give an example * on the real line * where these three sets are distinct.
\item
  Let A, B be two subsets of X; show that
\end{enumerate}

\[
\operatorname{Fr} (\mathbf {A} \cup \mathbf {B}) \subset \operatorname{Fr} (\mathbf {A}) \cup \operatorname{Fr} (\mathbf {B})
\]

and give an example * on the real line * where these sets are distinct. If \(\overline{\mathbf{A}} \cap \overline{\mathbf{B}} = \emptyset\) , show that \(\operatorname{Fr}(\mathbf{A} \cup \mathbf{B}) = \operatorname{Fr}(\mathbf{A}) \cup \operatorname{Fr}(\mathbf{B})\) .

\begin{enumerate}
\def\labelenumi{\alph{enumi})}
\setcounter{enumi}{2}
\tightlist
\item
  If \(\mathbf{A}\) and \(\mathbf{B}\) are open in \(\mathbf{X}\) , show that
\end{enumerate}

\[
\begin{array}{r l} & (\mathbf {A} \cap \operatorname{Fr} (\mathbf {B})) \cup (\mathbf {B} \cap \operatorname{Fr} (\mathbf {A})) \subset \operatorname{Fr} (\mathbf {A} \cap \mathbf {B}) \\ & \qquad \subset (\mathbf {A} \cap \operatorname{Fr} (\mathbf {B})) \cup (\mathbf {B} \cap \operatorname{Fr} (\mathbf {A})) \cup (\operatorname{Fr} (\mathbf {A}) \cap \operatorname{Fr} (\mathbf {B})) \end{array}
\]

and give an example * on the real line * where these three sets are distinct.

\begin{enumerate}
\def\labelenumi{\arabic{enumi})}
\setcounter{enumi}{5}
\item
  Show that a subset A of a topological space X meets each dense subset of X if and only if the interior of A is not empty.
\item
  Consider the following four properties of a given topological space \(\mathbf{X}\) :
\end{enumerate}

(D \(_{I}\) ) The topology of X has a countable base.

\((D_{II})\) X has a countable dense subset.

\((D_{III})\) Every subset of X, all of whose points are isolated, is countable.

\((D_{IV})\) Every set of mutually disjoint non-empty open subsets of X is countable.

Show that \((\mathbf{D}_{\mathbf{I}})\) implies \((\mathbf{D}_{\mathbf{II}})\) and \((\mathbf{D}_{\mathbf{III}})\) , and that each of \((\mathbf{D}_{\mathbf{II}})\) and \((\mathbf{D}_{\mathbf{III}})\) imply \((\mathbf{D}_{\mathbf{IV}})\) (*).

\begin{enumerate}
\def\labelenumi{\arabic{enumi})}
\setcounter{enumi}{7}
\item
  If a subset A of a topological space X has no isolated points, then the same is true of its closure \(\overline{A}\) in X.
\item
  Let X be a set and let \(M \to \overline{M}\) be a mapping of \(\mathfrak{P}(X)\) onto itself such that: 1) \(\overline{\varnothing} = \varnothing\) ; 2) for every \(M \subset X\) we have \(M \subset \overline{M}\) ; 3) for every \(M \subset X\) we have \(\overline{M} = \overline{M}\) ; 4) for all \(M \subset X\) and \(N \subset X\) we have
\end{enumerate}

\[
\overline {{\mathbf {M} \cup \mathbf {N}}} = \overline {{\mathbf {M}}} \cup \overline {{\mathbf {N}}}.
\]

Show that there is a unique topology on X such that \(\overline{M}\) is the closure of M with respect to this topology, for all \(M \subset X\) (define the topology by means of its closed sets).

\BourbakiExerciseGroup

\begin{enumerate}
\def\labelenumi{\arabic{enumi})}
\tightlist
\item
  Let \(f\) be a mapping of a topological space \(\mathbf{X}\) into a topological space \(\mathbf{X}'\) . Prove that the following statements are equivalent:
\end{enumerate}

\begin{enumerate}
\def\labelenumi{\alph{enumi})}
\item
  \(f\) is continuous on \(\mathbf{X}\) .
\item
  For every subset \(\mathbf{A}'\) of \(\mathbf{X}', f^1 (\mathring{\mathbf{A}}') \subset (f^1 (\mathbf{A}'))^0\) .
\item
  For every subset \(\mathbf{A}'\) of \(\mathbf{X}', \overline{f}^1 (\mathbf{A}') \subset \overline{f}^1 (\overline{\mathbf{A}}')\) .
\end{enumerate}

Show by an example that, if \(f\) is continuous, the sets \(\overline{f}^1 (\mathbf{A}')\) and \(\overline{f}^1 (\overline{\mathbf{A}}')\) can be distinct.

\begin{enumerate}
\def\labelenumi{\arabic{enumi})}
\setcounter{enumi}{1}
\item
  Let X, X' be two ordered sets, each topologized with the right topology (§ 1, Exercise 2). Show that a mapping \(f: \mathbf{X} \to \mathbf{X}'\) is continuous if and only if it is order-preserving.
\item
  Let X, X' be two topological spaces, f a bijection of X onto X'; then a necessary and sufficient condition for f to be a homeomorphism of X onto \(X'\) is that the topology of \(X'\) is the finest for which f is continuous.
\item
  On an ordered set X, the least upper bound of the left topology and the right topology (§ 1, Exercise 2) is the discrete topology; if X is a directed set (on the right or the left), the greatest lower bound of these two topologies is the coarsest topology on X.
\item
  On an ordered set X let \(\mathcal{C}_{0}(\mathbf{X})\) (resp. \(\mathcal{C}_{+}(\mathbf{X})\) , \(\mathcal{C}_{-}(\mathbf{X})\) ) denote the topology generated by the set of all open intervals (resp. the set of right half-open intervals, the set of left half-open intervals), bounded or not.
\end{enumerate}

\begin{enumerate}
\def\labelenumi{\alph{enumi})}
\item
  Show that if X is linearly ordered, then the open (resp. right half-open, left half-open) intervals form a base of \(\mathcal{C}_{0}(\mathbf{X})\) (resp. \(\mathcal{C}_{+}(\mathbf{X})\) , \(\mathcal{C}_{-}(\mathbf{X})\) ); and that the closed intervals are closed sets in each of these three topologies.
\item
  On a linearly ordered set X, show that the topology \(\mathcal{C}_{+}(\mathbf{X})\) is finer than the right topology (§ 1, Exercise 2), that the topology \(\mathcal{C}_{-}(\mathbf{X})\) is finer than the left topology, and that the topology \(\mathcal{C}_{0}(\mathbf{X})\) is the intersection of \(\mathcal{C}_{+}(\mathbf{X})\) and \(\mathcal{C}_{-}(\mathbf{X})\) .
\item
  Let X be a linearly ordered set. Show that a necessary and sufficient condition for \(\mathfrak{C}_{0}(\mathbf{X})\) to have a countable base is that there should be is a countable subset D of X such that, for each \(x \in X\) and each interval {]}a, b{[} containing x, there exist \(\alpha \in D\) and \(\beta \in D\) such that
\end{enumerate}

\[
a \leqslant \alpha <   x <   \beta \leqslant b
\]

{[}cf.~Chapter IV, § 2, Exercise 11 c){]}.

\begin{enumerate}
\def\labelenumi{\alph{enumi})}
\setcounter{enumi}{3}
\tightlist
\item
  Show that if X is well-ordered (Set Theory, Chapter III, § 2) then the topology \(\mathcal{G}_{+}(\mathbf{X})\) is the discrete topology, and the topologies \(\mathcal{G}_{0}(\mathbf{X})\) and \(\mathcal{G}_{-}(\mathbf{X})\) are identical.
\end{enumerate}

¶ 6) A topology on a set X is said to be quasi-maximal if it is maximal in the set of topologies in which X has no isolated points.

\begin{enumerate}
\def\labelenumi{\alph{enumi})}
\item
  Show that the following properties of a topology \(\mathcal{G}\) on X are equivalent: \(\alpha\) \(\mathcal{G}\) is quasi-maximal; \(\beta\) if X carries the topology \(\mathcal{G}\) , then X has no isolated point and every subset of X which has no isolated points is open.
\item
  Let X be a Kolmogoroff space (§ 1, Exercise 2) in which every non-empty open set is infinite. Show that there exists a quasi-maximal topology on X which is finer than the given topology (show that the set of topologies on X in which all the non-empty open sets are infinite is inductive, and use Exercise 2 d) of § 1).
\end{enumerate}

\begin{enumerate}
\def\labelenumi{\arabic{enumi})}
\setcounter{enumi}{6}
\tightlist
\item
  If X is any set, let \(\mathfrak{P}_{0}(\mathbf{X})\) denote the set of non-empty subsets of X.
\end{enumerate}

\begin{enumerate}
\def\labelenumi{\alph{enumi})}
\item
  Let X be a non-empty topological space. Let \(\mathcal{G}_{\Omega}\) (resp. \(\mathcal{G}_{\Phi}\) ) denote the coarsest topology on \(\mathfrak{P}_{0}(\mathbf{X})\) such that for every non-empty open (resp. closed) subset A of X, \(\mathfrak{P}_{0}(\mathbf{A})\) is open (resp. closed) in \(\mathfrak{P}_{0}(\mathbf{X})\) . Show that in general these two topologies are not comparable, and that the mapping \(x \to \{x\}\) is a homeomorphism of X onto a subspace of \(\mathfrak{P}_{0}(\mathbf{X})\) if this latter set carries either of the topologies \(\mathcal{G}_{\Omega}, \mathcal{G}_{\Phi}\) .
\item
  Let D be a dense subset of X. Show that the set \(\mathfrak{E}(D)\) of finite non-empty subsets of D is dense in \(\mathfrak{P}_{0}(X)\) for either of the topologies \(C_{\Omega}, C_{\Phi}\) .
\item
  If \(\mathfrak{P}_{0}(\mathbf{X})\) is ordered by inclusion, show that the topology \(\mathcal{C}_{\Omega}\) is coarser than the left topology on \(\mathfrak{P}_{0}(\mathbf{X})\) ( \(\S\) 1, Exercise 2); that \(\mathbf{A} \in \mathfrak{P}_{0}(\mathbf{X})\) is isolated in the topology \(\mathcal{C}_{\Omega}\) if and only if \(\mathbf{A} = \{x\}\) , where \(x\) is an isolated point of \(\mathbf{X}\) ; and that \(\mathbf{A} \in \mathfrak{P}_{0}(\mathbf{X})\) is isolated in the topology \(\mathcal{C}_{\Phi}\) if and only if \(\mathbf{X}\) is a finite discrete space and \(\mathbf{A} = \mathbf{X}\) .
\item
  Prove that if \(\mathfrak{P}_0(\mathbf{X})\) carries the topology \(\mathcal{C}_{\Omega}\) (resp. \(\mathcal{C}_{\Phi}\) ) and
\end{enumerate}

\[
\mathfrak {P} _ {\mathbf {0}} (\mathbf {X}) \times \mathfrak {P} _ {\mathbf {0}} (\mathbf {X})
\]

carries the topology which is the product of \(\mathcal{C}_{\Omega}\) (resp. \(\mathcal{C}_{\Phi}\) ) with itself, then the mapping \((\mathbf{M}, \mathbf{N}) \to \mathbf{M} \cup \mathbf{N}\) of \(\mathfrak{P}_0(\mathbf{X}) \times \mathfrak{P}_0(\mathbf{X})\) into \(\mathfrak{P}_0(\mathbf{X})\) is continuous.

\begin{enumerate}
\def\labelenumi{\alph{enumi})}
\setcounter{enumi}{4}
\tightlist
\item
  Show that if \(\mathfrak{P}_{0}(\mathbf{X})\) carries the topology \(\mathfrak{T}_{\Phi}\) , then the mapping \(\mathbf{M} \to \overline{\mathbf{M}}\) is continuous.
\end{enumerate}

\(f)\) Let \(\mathbf{X}\) , \(\mathbf{Y}\) be two topological spaces, \(f: \mathbf{X} \to \mathbf{Y}\) a continuous map. Show that if both \(\mathfrak{P}_{0}(\mathbf{X})\) and \(\mathfrak{P}_{0}(\mathbf{Y})\) carry the topology \(\mathcal{G}_{\Omega}\) , and if both carry the topology \(\mathcal{G}_{\Phi}\) , then the mapping \(\mathbf{M} \to f(\mathbf{M})\) of \(\mathfrak{P}_{0}(\mathbf{X})\) into \(\mathfrak{P}_{0}(\mathbf{Y})\) is continuous.

\begin{enumerate}
\def\labelenumi{\alph{enumi})}
\setcounter{enumi}{6}
\tightlist
\item
  Let X, Z be two topological spaces. Show that a mapping \(f: Z \to \mathfrak{P}_{0}(X)\) is continuous when \(\mathfrak{P}_{0}(X)\) carries the topology \(C_{\Omega}\) (resp. \(C_{\Phi}\) ) if and only if the set of all \(z \in Z\) such that \(f(z) \cap A \neq \emptyset\) is closed (resp. open) in Z for each closed (resp. open) subset A of X.
\end{enumerate}

\begin{enumerate}
\def\labelenumi{\arabic{enumi})}
\setcounter{enumi}{7}
\tightlist
\item
  Let X be a set, and c a cardinal number which is either infinite or equal to 1. The set of subsets of X consisting of X and those M ⊆ X such that card(M) \textless{} c is the set of closed subsets of X for a topology \(\mathfrak{G}_{\mathfrak{c}}\) ; \(\mathfrak{G}_{\mathfrak{1}}\) is the coarsest topology on X and \(\mathfrak{G}_{\mathfrak{c}}\) is the discrete topology if c \textgreater{} card (X). Every permutation of X is an automorphism of the topological space obtained by giving X a topology \(\mathfrak{G}_{\mathfrak{c}}\) .
\end{enumerate}

Show that conversely any topology \(\mathcal{C}\) on \(\mathbf{X}\) , which has the property that every permutation of \(\mathbf{X}\) is continuous with respect to \(\mathcal{C}\) , is necessarily one of the topologies \(\mathcal{C}_{\mathfrak{c}}\) . {[}Let \(\mathfrak{c}\) be the smallest cardinal such that card \((\mathbf{F}) < \mathfrak{c}\) for each closed subset \(\mathbf{F} \neq \mathbf{X}\) in the topology \(\mathcal{C}\) ; observe that when \(\mathbf{X}\) is infinite and \(\mathfrak{c} >\mathrm{card}(\mathbf{X})\) , \(\mathcal{C}\) is necessarily the discrete topology.{]}

\begin{enumerate}
\def\labelenumi{\arabic{enumi})}
\setcounter{enumi}{8}
\tightlist
\item
  \begin{enumerate}
  \def\labelenumii{\alph{enumii})}
  \tightlist
  \item
    With the hypotheses of Proposition 4 of no. 3, show that the coarsest topology on X for which the \(f_{t}\) are continuous is also the finest topology \(\mathfrak{G}\) on X which has the following property: every mapping \(g\) of a topological space Z into X such that the \(f_{t} \circ g\) are continuous is continuous (with respect to \(\mathfrak{G}\) ).
  \end{enumerate}
\end{enumerate}

\begin{enumerate}
\def\labelenumi{\alph{enumi})}
\setcounter{enumi}{1}
\tightlist
\item
  With the hypotheses of Proposition 6 of no. 4, show that the finest topology on X for which the \(f_{i}\) are continuous is also the coarsest topology \(\mathcal{G}\) on X which has the following property: every mapping \(g\) of X into a topological space Z such that the \(g \circ f_{i}\) are continuous is continuous (with respect to \(\mathcal{G}\) ).
\end{enumerate}

\begin{enumerate}
\def\labelenumi{\arabic{enumi}.}
\setcounter{enumi}{9}
\tightlist
\item
  On a set X, let M be a directed set of topologies without isolated points. Show that the least upper bound of M in the set of all topologies on X is a topology without isolated points.
\end{enumerate}

¶ 11. A topological space is said to be solvable if it is the disjoint union of two dense subsets.

\begin{enumerate}
\def\labelenumi{\alph{enumi})}
\item
  If a topological space X is solvable, show that X has no isolated points and that every open subspace of X is solvable. Show also conversely that if X is a topological space and B is a set of non-empty solvable open subspaces of X such that every non-empty open subset of X contains a set belonging to B, then X is solvable (consider a maximal set of mutually disjoint open sets belonging to B).
\item
  If a Kolmogoroff space X is solvable, then every non-empty open subset of X is infinite {[}use Exercise 2d) of §1{]}.
\item
  Let X be a topological space such that 1) the smallest cardinal b of a non-empty open subset of X is infinite, and 2) there is a base B of the topology of X such that card (B) ≤ b. Show that X is solvable. (Use transfinite induction to construct two disjoint dense subsets of X). By the method of Set Theory, Chapter III, § 6, Exercise 24 a). In particular, the rational line is solvable.
\item
  A topological space X which has no isolated points is said to be isodyne if card (U) = card (X) for every non-empty open subset of X. Show that in a space X with no isolated points, every non-empty open set in X contains a non-empty open subspace which is isodyne. Deduce that if every isodyne open subspace of X is soluble, then X is soluble.
\end{enumerate}

\BourbakiExerciseGroup

\begin{enumerate}
\def\labelenumi{\arabic{enumi})}
\tightlist
\item
  Let A, B be two subsets of a topological space X such that \(A \supset B\) . Show that:
\end{enumerate}

\begin{enumerate}
\def\labelenumi{\alph{enumi})}
\item
  The interior of B with respect to X is contained in the interior of B with respect to the subspace A of X. Give an example where these two sets are distinct.
\item
  The frontier of B with respect to A is contained in the intersection of A with the frontier of B with respect to X. Give an example where these two sets are distinct.
\end{enumerate}

\begin{enumerate}
\def\labelenumi{\arabic{enumi})}
\setcounter{enumi}{1}
\tightlist
\item
  Let A and B be any two subsets of a topological space X.
\end{enumerate}

\begin{enumerate}
\def\labelenumi{\alph{enumi})}
\item
  Show that the intersection of A and the interior of B with respect to X is contained in the interior of \(B \cap A\) with respect to A. Give an example where these two sets are distinct.
\item
  Show that the intersection of A and the closure of B with respect to X contains the closure of B ∩ A with respect to A. Give an example where the two sets are distinct.
\end{enumerate}

\begin{enumerate}
\def\labelenumi{\arabic{enumi})}
\setcounter{enumi}{2}
\item
  A subspace \(\mathbf{A}\) of a topological space \(\mathbf{X}\) is discrete if and only if every point of \(\mathbf{A}\) is isolated.
\item
  Let Y, Z be two subspaces of a topological space X such that \(X = Y \cup Z\) and such that \(\overline{A} \cap B = \overline{B} \cap A = \varnothing\) , where \(A = Y \cap CZ\) and \(B = Z \cap CY\) .
\end{enumerate}

\begin{enumerate}
\def\labelenumi{\alph{enumi})}
\item
  Show that if M is any subset of X then the closure of M in X is the union of the closure of \(M \cap Y\) in Y and the closure of \(M \cap Z\) in Z. Deduce that if \(M \cap Y\) is closed (resp. open) in Y and \(M \cap Z\) is closed (resp. open) in Z, then M is closed (resp. open) in X.
\item
  Let \(f\) be a mapping of \(\mathbf{X}\) into a topological space \(\mathbf{X}'\) . Show that if \(f|_{\mathbf{Y}}\) and \(f|_{\mathbf{Z}}\) are continuous then so is \(f\) .
\end{enumerate}

\begin{enumerate}
\def\labelenumi{\arabic{enumi})}
\setcounter{enumi}{4}
\item
  Let Y, Z be two subspaces of a topological space X such that \(X = Y \cup Z\) . Show that if a subset M of \(Y \cap Z\) is open (resp. closed) in both Y and Z, then M is open (resp. closed) in X.
\item
  Let A be a locally closed subspace of a topological space X. Show that the set of open subsets U of X such that \(A \subset U\) and A is closed in U has a greatest element, namely the complement in X of the frontier of A with respect to \(\overline{A}\) .
\item
  Show by examples that the union of two locally closed subsets and the complement of a locally closed subset need not be locally closed.
\item
  Let A be a subset of an ordered set X. Show that the topology induced on A by the right (resp. left) topology of X (§ 1, Exercise 2) is the right (resp. left) topology of A with respect to the ordering of A induced by that of X.
\item
  Let \(\mathbf{A}\) be a subset of a linearly ordered set \(\mathbf{X}\) .
\end{enumerate}

\begin{enumerate}
\def\labelenumi{\alph{enumi})}
\item
  Show that the topology induced on A by the topology \(\mathcal{G}_{-}(\mathbf{X})\) {[}resp. \(\mathcal{G}_{+}(\mathbf{X}),\mathcal{G}_{0}(\mathbf{X})]\) (§ 2, Exercise 5) is finer than the topology \(\mathcal{G}_{-}(\mathbf{A})\) {[}resp. \(\mathcal{G}_{+}(\mathbf{A}),\mathcal{G}_{0}(\mathbf{A})]\) , the ordering of A being that induced by the ordering of X. If A is an interval of X, then \(\mathcal{G}_{-}(\mathbf{A})\) {[}resp. \(\mathcal{G}_{+}(\mathbf{A}),\mathcal{G}_{0}(\mathbf{A})]\) is the topology induced on A by \(\mathcal{G}_{-}(\mathbf{X})\) {[}resp. \(\mathcal{G}_{+}(\mathbf{X}),\mathcal{G}_{0}(\mathbf{X})]\) .
\item
  Take X to be the product \(Q \times Z\) , ordered by the lexicographic ordering (Set Theory, Chapter III, § 2, no. 6) and take A to be the subset \(Q \times \{0\}\) . Show that the topology \(\mathcal{G}_{-}(A)\) {[}resp. \(\mathcal{G}_{+}(A)\) , \(\mathcal{G}_{0}(A)\) {]} is strictly coarser than the topology induced on A by \(\mathcal{G}_{-}(X)\) {[}resp. \(\mathcal{G}_{+}(X)\) , \(\mathcal{G}_{0}(X)\) {]}.
\end{enumerate}

\begin{enumerate}
\def\labelenumi{\arabic{enumi})}
\setcounter{enumi}{9}
\item
  Let A be a non-empty subspace of a non-empty topological space X. Show that the topology induced on \(\mathfrak{P}_{0}(\mathbf{A})\) by the topology \(\mathcal{G}_{\Omega}\) (resp. \(\mathcal{G}_{\Phi}\) ) on \(\mathfrak{P}_{0}(\mathbf{X})\) (§ 2, Exercise 7) is the topology \(\mathcal{G}_{\Omega}\) (resp. \(\mathcal{G}_{\Phi}\) ) on \(\mathfrak{P}_{0}(\mathbf{A})\) .
\item
  Let X be a topological space, \(\mathcal{G}\) the topology of X, A a dense subspace of X. Show that the set of topologies on X which are finer than \(\mathcal{G}\) , in which A is dense in X, and which induce the same topology as \(\mathcal{G}\) on A, has at least one maximal element; such a maximal element is called an A-maximal topology. Show that a topology \(\mathcal{G}_0\) on X is A-maximal if and only if every subset M of X, such that A∩M is dense in M (with respect to \(\mathcal{G}_0\) ) and open in A, is open in the topology \(\mathcal{G}_0\) , and that the subspace \(\mathbb{C}\) A is then discrete in the topology induced by \(\mathcal{G}_0\) , and is closed in X. Show also that if \(\mathcal{G}_0\) is A-maximal and if the topology induced by \(\mathcal{G}_0\) on A is quasi-maximal (§ 2, Exercise 6) then \(\mathcal{G}_0\) is quasi-maximal.
\end{enumerate}

* 12) Let X be the subspace {[}0, 1{]} of the real line R, and let S be the equivalence relation on X whose equivalence classes are \{0, 1\} and the sets \{x\} for 0 \textless{} x \textless{} 1. Show that X has no continuous section with respect to S. *

\begin{enumerate}
\def\labelenumi{\arabic{enumi})}
\setcounter{enumi}{12}
\item
  Let X be a topological space, and R and S two equivalence relations on X such that R implies S. Show that if there is a continuous section of X with respect to R and a continuous section of X/R with respect to S/R, then there is a continuous section of X with respect to S.
\item
  Let X be a topological space which is the sum of two subspaces \(X_{1}\) and \(X_{2}\) , and let X/R be the space obtained by pasting together \(X_{1}\) and \(X_{2}\) along a subspace \(A_{1}\) of \(X_{1}\) and a subspace \(A_{2}\) of \(X_{2}\) by means of a homeomorphism h of \(A_{1}\) onto \(A_{2}\) . Show that the canonical images of \(X_{1}\) and \(X_{2}\) in X/R are homeomorphic to \(X_{1}\) and \(X_{2}\) respectively.
\end{enumerate}

* 15) Consider the following three subspaces of \(\mathbf{R}: \mathbf{X}_1 = ] - \mathbf{1}, \mathbf{1}[, \mathbf{X}_2 = ] - \mathbf{2}, -\mathbf{1}[, \mathbf{X}_3 = ]\mathbf{1}, \mathbf{2}[.\) Let \(\mathbf{X}\) be their union; \(\mathbf{X}\) is then the sum of these three subspaces. Let \(A_{2}\) (resp. \(A_{3}\) ) be the set of irrational numbers contained in \(X_{2}\) (resp. \(X_{3}\) ); let \(B_{1}\) (resp. \(B_{2}\) ) be the set of rational numbers contained in {]}-1, 0{[} (resp. \(X_{2}\) ); and let \(C_{1}\) (resp. \(C_{3}\) ) be the set of rational numbers contained in {]}0, 1{[} (resp. \(X_{3}\) ). Let X/R be the space obtained by pasting together the \(X_{i}\) (i) along \(A_{2}\) and \(A_{3}\) by means of the homeomorphism \(x \to x + 3\) ; (ii) along \(B_{1}\) and \(B_{2}\) by means of the homeomorphism \(x \to x - 1\) ; (iii) along \(C_{1}\) and \(C_{3}\) by means of the homeomorphism \(x \to x + 1\) . Show that the canonical image of \(X_{1}\) in X/R is not homeomorphic to \(X_{1}\) . If \(Y = X_{1} \cup B_{2} \cup C_{3}\) is the saturation of \(X_{1}\) with respect to R, show that the quotient space \(Y/R_{Y}\) is not homeomorphic to the canonical image of Y in X/R. *

\begin{enumerate}
\def\labelenumi{\arabic{enumi})}
\setcounter{enumi}{15}
\tightlist
\item
  Let X be a non-empty topological space and R an equivalence relation on X. Show that if X/R is regarded as a subset of \(\mathfrak{P}_{0}(X)\) , then the quotient topology on X/R is coarser than the topologies induced on X/R by the topologies \(C_{\Omega}\) and \(C_{\Phi}\) on \(\mathfrak{P}_{0}(X)\) (§ 2, Exercise 7).
\end{enumerate}

\BourbakiExerciseGroup

\begin{enumerate}
\def\labelenumi{\arabic{enumi})}
\item
  Let X, Y be two topological spaces, A a subset of X and B a subset of Y. Show that \(\mathrm{Fr}(\mathbf{A} \times \mathbf{B}) = (\mathrm{Fr}(\mathbf{A}) \times \overline{\mathbf{B}}) \cup (\overline{\mathbf{A}} \times \mathrm{Fr}(\mathbf{B}))\) .
\item
  Let \(\mathbf{X}\) , \(\mathbf{Y}\) be two topological spaces, A a closed subset of \(\mathbf{X} \times \mathbf{Y}\) . If \(\mathfrak{B}(x)\) is the set of all neighbourhoods of a point \(x \in \mathbf{X}\) , show that
\end{enumerate}

\[
\bigcap_ {\mathbf {V} \in \mathfrak {B} (x)} \overline {{\mathrm{A} (\mathbf {V})}} = \mathrm{A} (x).
\]

Give an example of an open subset B of a product \(X \times Y\) such that

\[
\bigcap_ {\mathbf {V} \in \mathfrak {B} (x)} \overline {{\mathbf {B} (\mathbf {V})}}
\]

is different from \(\mathbf{B}_{(x)}\) .

\begin{enumerate}
\def\labelenumi{\arabic{enumi})}
\setcounter{enumi}{2}
\tightlist
\item
  Let \((\mathbf{X}_t)_{i\in I}\) be an infinite family of topological spaces such that each \(\mathbf{X}_t\) contains (at least) two distinct points \(a_i, b_i\) such that there is a neighbourhood of \(b_i\) which does not contain \(a_i\) .
\end{enumerate}

\begin{enumerate}
\def\labelenumi{\alph{enumi})}
\tightlist
\item
  For each index \(x \in I\) , let \(c_{x}\) be the point of the product space
\end{enumerate}

\[
\mathbf {X} = \prod_ {k + 1} \mathbf {X} _ {i}
\]

for which \(\mathrm{pr}_x c_x = b_x\) and \(\mathrm{pr}_t c_x = a_t\) whenever \(i \neq x\) . Show that every point of the set C of points \(c_x\) is isolated.

\begin{enumerate}
\def\labelenumi{\alph{enumi})}
\setcounter{enumi}{1}
\item
  Deduce from a) that the topology of X has a countable base if and only if I is countable and the topology of each \(X_{i}\) has a countable base (use Exercise 7 of §1).
\item
  Show that if the index set I is not countable then the point \(b = (b_{i})\) of X has no countable fundamental system of neighbourhoods.
\end{enumerate}

\begin{enumerate}
\def\labelenumi{\arabic{enumi})}
\setcounter{enumi}{3}
\tightlist
\item
  Let K be the discrete space consisting of the two numbers 0,1; let A be an infinite set, and let X be the product space \(\mathbf{K}^{\mathbf{A}}\) . Let
\end{enumerate}

\[
\mathbf {V} = \prod_ {\alpha \in \mathbf {A}} \mathbf {V} _ {\alpha}
\]

be an elementary set in X; if h is the number of indices \(\alpha\) such that \(V_{\alpha} \neq K\) , put \(\mu(V) = 2^{-h}\) .

\begin{enumerate}
\def\labelenumi{\alph{enumi})}
\tightlist
\item
  Show that if \(\mathbf{U}_1, \ldots, \mathbf{U}_n\) are pairwise disjoint elementary sets, then
\end{enumerate}

\[
\sum_ {k = 1} ^ {n} \mu (\mathrm{U} _ {k}) \leqslant \mathrm{I}.
\]

{[}Write each \(U_{k}\) in the form \(W_{k} \times K^{B}\) , where B (independent of k) is the complement of a finite subset of A, and \(W_{k}\) is a finite set; count up the number of points in \(W_{k}\) {]}.

\begin{enumerate}
\def\labelenumi{\alph{enumi})}
\setcounter{enumi}{1}
\item
  Deduce from a) that X satisfies the condition \((\mathbf{D}_{\mathbf{IV}})\) of Exercise 7 of §1. {[}If there were an uncountable family \((\mathbf{U}_{\lambda})\) of pairwise disjoint elementary sets, there would exist an integer p \textgreater{} 0 such that \(\mu(\mathbf{U}_{\lambda}) \geqslant 1/p\) for an infinite number of indices \(\lambda\) {]}.
\item
  Show that if A is uncountable, then X does not satisfy the condition \((\mathbf{D}_{\mathrm{III}})\) of Exercise 7 of §1 (cf.~Exercise 3).
\end{enumerate}

\begin{enumerate}
\def\labelenumi{\arabic{enumi})}
\setcounter{enumi}{4}
\tightlist
\item
  Let K be the discrete space consisting of the two numbers 0, 1; let A be an infinite set and let X denote the product space \(\mathbf{K}^{\mathfrak{P}(\mathbf{A})}\) . a) Let J be any finite subset of A; suppose J has p elements. For each subset H of J, let \(M_{H}\) be the set of subsets L of A such that \(L \cap J = H\) ; the \(M_{H}\) form a partition \(\tilde{\omega}_{J}\) of \(\mathfrak{P}(A)\) into \(2^{p}\) subsets. Let \(F_{J}\) be the subset of X consisting of those points \((x_{L})_{L \subset A}\) such that \(x_{L} = x_{M}\) whenever L and M belong to the same set of the partition \(\tilde{\omega}_{J}\) . \(F_{J}\) is a finite set, with \(2^{2^{p}}\) elements. If F is the union of the \(F_{J}\) as J runs through all finite subsets of A, show that F is equipotent to A.
\end{enumerate}

\begin{enumerate}
\def\labelenumi{\alph{enumi})}
\setcounter{enumi}{1}
\tightlist
\item
  Show that F is dense in X. Taking A = N, deduce that X satisfies condition \((D_{\mathrm{II}})\) of §1, Exercise 7, but does not satisfy condition \((D_{\mathrm{III}})\) (Exercise 4).
\end{enumerate}

\begin{enumerate}
\def\labelenumi{\arabic{enumi})}
\setcounter{enumi}{5}
\item
  Let X, Y, Z be three topological spaces, A an open subset of \(X \times Y\) and B an open subset of \(Y \times Z\) . Show that \(B \circ A\) is an open subset of \(X \times Z\) .
\item
  Let X be the sum of a family \((X_{\lambda})\) of topological spaces, Y the sum of a family \((Y_{\mu})\) of topological spaces. Show that the product space \(X \times Y\) is the sum of the \(X_{\lambda} \times Y_{\mu}\) .
\item
  Let \((\mathbf{G}_t)_{t \in \mathbf{I}}\) be a family of graphs of equivalence relations (Set Theory, Chapter II, § 6, no. 1) on a topological space \(\mathbf{X}\) , and let \(\mathbf{R}_t\) be the relation \((x, y) \in \mathbf{G}_t\) ; then
\end{enumerate}

\[
\mathbf {G} = \bigcap_ {\iota \in \mathbf {I}} \mathbf {G} _ {\iota}
\]

is the graph of an equivalence relation R on X. Show that there is a canonical continuous injection of the quotient space X/R into the product space \(\prod_{i}\left(X/R_{i}\right)\) .

\begin{enumerate}
\def\labelenumi{\arabic{enumi})}
\setcounter{enumi}{8}
\tightlist
\item
  Let \((\mathbf{X}_i)_{i\in I}\) be an infinite family of topological spaces. Study the topology on the product set
\end{enumerate}

\[
\mathbf {X} = \prod_ {i \in I} \mathbf {X} _ {i} \text {   generated   by   subsets   of   the   form   } \prod_ {i \in I} \mathbf {U} _ {i},
\]

where \(U_{t}\) is an open subset of \(X_{t}\) for each \(t \in I\) and the set of indices t for which \(U_{t} \neq X_{t}\) has cardinal less than c, where c is a given infinite cardinal. Consider what becomes of the propositions of §4 for this topology.

\begin{enumerate}
\def\labelenumi{\arabic{enumi})}
\setcounter{enumi}{9}
\tightlist
\item
  \begin{enumerate}
  \def\labelenumii{\alph{enumii})}
  \tightlist
  \item
    Let \((\mathbf{X}_{\alpha}, f_{\alpha\beta})\) be an inverse system of topological spaces, and let X be the inverse limit. Let \(f_{\alpha}\) be the canonical mapping \(X \to X_{\alpha}\) . For each \(\alpha\) , let \(Y_{\alpha}\) be a subspace of \(X_{\alpha}\) containing \(f_{\alpha}(X)\) , such that the \(Y_{\alpha}\) form an inverse system of subspaces of the \(X_{\alpha}\) . Show that \(\varprojlim Y_{\alpha}\) can be canonically identified with X.
  \end{enumerate}
\end{enumerate}

\begin{enumerate}
\def\labelenumi{\alph{enumi})}
\setcounter{enumi}{1}
\tightlist
\item
  Let \((\mathbf{X}_{\alpha}^{\prime}, f_{\alpha\beta}^{\prime})\) be another inverse system of topological spaces with the same directed index set. For each \(\alpha\) , let \(u_{\alpha}: X_{\alpha} \to X_{\alpha}^{\prime}\) be a continuous mapping such that the \(u_{\alpha}\) form an inverse system of mappings. The \(u_{\alpha}(\mathbf{X}_{\alpha})\) form an inverse system of subspaces of the \(X_{\alpha}^{\prime}\) ; show that if \(u = \varprojlim u_{\alpha}\) and if the \(f_{\alpha}\) are surjective, then \(u(\mathbf{X})\) is dense in the space \(\varprojlim u_{\alpha}(\mathbf{X}_{\alpha})\) . Does this result still hold good if the \(f_{\alpha}\) are no longer assumed to be surjective? (Cf. Set Theory, Chapter III, §1, Exercise 32.)
\end{enumerate}

\BourbakiExerciseGroup

\begin{enumerate}
\def\labelenumi{\arabic{enumi})}
\tightlist
\item
  \begin{enumerate}
  \def\labelenumii{\alph{enumii})}
  \tightlist
  \item
    Show that an equivalence relation R on a topological space X is open if and only if, for each subset A of X, the saturation of the interior of A is contained in the interior of the saturation of A. Give an example where R is open and these two sets can be distinct.
  \end{enumerate}
\end{enumerate}

\begin{enumerate}
\def\labelenumi{\alph{enumi})}
\setcounter{enumi}{1}
\tightlist
\item
  Show that an equivalence relation R on a topological space X is closed if and only if, for each subset A of X, the saturation of the closure of A contains the closure of the saturation of A. Give an example where R is closed and these two sets can be distinct.
\end{enumerate}

* 2) Show that the equivalence relation \(S: x \equiv y \pmod{1}\) on \(R\) is not closed. If \(A = [0,1]\) show that the induced relation \(S_A\) on \(A\) is closed but not open. *

\begin{enumerate}
\def\labelenumi{\arabic{enumi})}
\setcounter{enumi}{2}
\tightlist
\item
  Let \(\Gamma\) be a finite group of homeomorphisms of a topological space \(X\) , and let \(R\) be the equivalence relation: ``there exists \(\sigma \in \Gamma\) such that \(y = \sigma(x)\)''. Show that the relation \(R\) is both open and closed.
\end{enumerate}

* 4) Let \(\Gamma\) be the group of homeomorphisms of the quotient space \(\mathbf{T}\) of \(\mathbf{R}\) consisting of the identity and the homeomorphism induced on \(\mathbf{T}\) by the homeomorphism \(x \to 1/2 + x\) of \(\mathbf{R}\) onto itself. Let \(\mathbf{S}\) be the equivalence relation ``there exists \(\sigma \in \Gamma\) such that \(y = \sigma(x)\) '' on \(\mathbf{T}\) ; \(\mathbf{S}\) is both open and closed (Exercise 3). Show that \(\mathbf{T}\) has no continuous section with respect to \(\mathbf{S}\) , although both \(\mathbf{T}\) and \(\mathbf{T}/\mathbf{S}\) are compact, connected and locally connected. Let \(\mathbf{A}\) be the canonical image in \(\mathbf{T}\) of the interval \([0,1/2]\) of \(\mathbf{R}\) . Show that \(\mathbf{A}\) is locally closed in \(\mathbf{T}\) but that the quotient space \(\mathbf{A}/\mathbf{S}_{\mathbf{A}}\) is not homeomorphic to the canonical image of \(\mathbf{A}\) in \(\mathbf{T}/\mathbf{S}\) . *

\begin{enumerate}
\def\labelenumi{\arabic{enumi})}
\setcounter{enumi}{4}
\tightlist
\item
  Let X be a topological space and let R, S be two equivalence relations on X such that R implies S.
\end{enumerate}

\begin{enumerate}
\def\labelenumi{\alph{enumi})}
\tightlist
\item
  Show that if S is open (resp. closed) then S/R is an open (resp. closed) equivalence relation on X/R; show also that the converse is not necessarily true. Show that S can be open (resp. closed) but not R.
\item
  Suppose that R is open (resp. closed). Show that S/R is open (resp. closed) if and only if S is open (resp. closed).
\end{enumerate}

\begin{enumerate}
\def\labelenumi{\arabic{enumi})}
\setcounter{enumi}{5}
\item
  Let S be the equivalence relation on the rational line Q which is obtained by identifying all the points of Z. Show that S is closed, and that if U is the relation of equality on Q, then the canonical bijection of \((\mathbf{Q} \times \mathbf{Q}) / (\mathbf{U} \times \mathbf{S})\) onto \(\mathbf{Q} \times (\mathbf{Q}/\mathbf{S})\) is not a homeomorphism.
\item
  Let X and Y be two non-empty topological spaces and \(f: X \to Y\) a surjective map. Show that f is closed if and only if the mapping \(y \to (\overline{f} y)\) of Y into \(\mathfrak{P}_{0}(X)\) is continuous when \(\mathfrak{P}_{0}(X)\) carries the topology \(C_{\Omega}\) (resp. \(C_{\Phi}\) ) (§ 2, Exercise 7).
\end{enumerate}

Deduce that the topology induced on a quotient set X/R of a non-empty topological space X by the topology \(\mathcal{G}_{\Omega}\) (resp. \(\mathcal{T}_{\Phi}\) ) coincides with the quotient of the topology of X by R if and only if R is closed (resp. open) (cf.~§ 3, Exercise 16).

\BourbakiExerciseGroup

\begin{enumerate}
\def\labelenumi{\arabic{enumi})}
\item
  Find all possible filters on a finite set.
\item
  If the intersection of all the sets of a filter F on an infinite set X is empty, show that F is finer than the filter of complements of finite subsets of X.
\item
  The intersection of two filters \(\mathfrak{F}\) and \(\mathfrak{G}\) on a set \(X\) is the set of all \(M \cup N\) , where \(M \in \mathfrak{F}\) and \(N \in \mathfrak{G}\) .
\item
  Let X be an infinite set. Show that the filter of complements of finite subsets of X is the intersection of the elementary filters associated with infinite sequences in X all of whose terms are distinct.
\item
  Let \(\mathfrak{F}\) and \(\mathfrak{G}\) be two filters on a set X. Show that if \(\mathfrak{F}\) and \(\mathfrak{G}\) have a least upper bound in the set of all filters on X, then this least upper bound is the set of all \(M \cap N\) , where \(M \in \mathfrak{F}\) and \(N \in \mathfrak{G}\) .
\item
  If to each filter on a set X we make correspond the associated topology (no. 5) on the set \(X' = X \cup \{\omega\}\) , then we have the following propositions:
\end{enumerate}

\begin{enumerate}
\def\labelenumi{\alph{enumi})}
\item
  If a topology \(\mathfrak{C}\) on \(\mathbf{X}'\) is finer than the topology associated with a filter \(\mathfrak{F}\) on \(\mathbf{X}\) , then \(\mathfrak{C}\) is either the discrete topology or the topology associated with a filter finer than \(\mathfrak{F}\) . Converse?
\item
  The greatest lower bound of the topologies associated with the filters of a set \(\Phi\) of filters on X is the topology associated with the intersection of the filters belonging to \(\Phi\) .
\item
  Let \(\mathfrak{G}\) be a set of subsets of X, and let \(\mathfrak{G}'\) be the set of subsets of X' of the form \(\mathbf{M} \cup \{\omega\}\) where \(\mathbf{M} \in \mathfrak{G}\) . Let \(\mathfrak{G}\) denote \(\mathfrak{G} \cup \mathfrak{P}(\mathbf{X})\) . If \(\mathfrak{G}\) generates a filter \(\mathfrak{F}\) , show that \(\mathfrak{G}\) is a subbase of the topology on X' associated with \(\mathfrak{F}\) . What topology does \(\mathfrak{G}\) generate if \(\mathfrak{G}\) is not a filter subbase?
\item
  \(\mathfrak{F}\) is an ultrafilter on \(\mathbf{X}\) if and only if the associated topology on \(\mathbf{X}'\) is \(\mathbf{X}\) -maximal ( \(\S 3\) , Exercise 11).
\end{enumerate}

\begin{enumerate}
\def\labelenumi{\arabic{enumi})}
\setcounter{enumi}{6}
\item
  In a topological space X, the intersection of the neighbourhood filters of all the points of a non-empty subset A of X is the neighbourhood filter of A in X.
\item
  \begin{enumerate}
  \def\labelenumii{\alph{enumii})}
  \tightlist
  \item
    Let \(\Phi\) be a set of topologies on a set \(X\) . Show that the neighbourhood filter of a point \(x \in X\) with respect to the intersection of the topologies of \(\Phi\) is coarser than the intersection of the neighbourhood filters of \(x\) with respect to the topologies belonging to \(\Phi\) .
  \end{enumerate}
\end{enumerate}

\begin{enumerate}
\def\labelenumi{\alph{enumi})}
\setcounter{enumi}{1}
\tightlist
\item
  Let \(\mathfrak{G}_1\) denote the topology \(\mathfrak{G}_0(\mathbf{Q}^2)\) (§ 2, Exercise 7) where the set \(\mathbf{Q}^2\) is linearly ordered by the lexicographic ordering (Set Theory,
\end{enumerate}

Chapter III, § 2, no. 6). Let \(\mathcal{G}_2\) denote the topology on \(\mathbf{Q}^2\) obtained by transporting \(\mathcal{G}_1\) by means of the symmetry \((\xi, \eta) \to (\eta, \xi)\) . Show that if \(\mathcal{G}\) is the intersection of the topologies \(\mathcal{G}_1, \mathcal{G}_2\) , then the neighbourhood filter of a point of \(\mathbf{Q}^2\) with respect to \(\mathcal{G}\) is strictly coarser than the intersection of the neighbourhood filters of the point with respect to \(\mathcal{G}_1\) and \(\mathcal{G}_2\) .

\begin{enumerate}
\def\labelenumi{\arabic{enumi})}
\setcounter{enumi}{8}
\tightlist
\item
  \begin{enumerate}
  \def\labelenumii{\alph{enumii})}
  \tightlist
  \item
    Show that every ultrafilter which is finer than the intersection of a finite number of filters is finer than at least one of them (use Proposition 5 of no. 4).
  \end{enumerate}
\end{enumerate}

\begin{enumerate}
\def\labelenumi{\alph{enumi})}
\setcounter{enumi}{1}
\tightlist
\item
  Give an example of an ultrafilter which is finer than the intersection of an infinite family of ultrafilters, but which is not equal to any of the ultrafilters of this family (consider the family of ultrafilters on an infinite set X each of which has a base consisting of a single point).
\end{enumerate}

\begin{enumerate}
\def\labelenumi{\arabic{enumi})}
\setcounter{enumi}{9}
\item
  Show that the intersection of all the sets of an ultrafilter contains at most one point; if one point, then the ultrafilter consists of all the sets containing this point (use Proposition 5 of no. 4).
\item
  Show that if a subset A of a set X does not belong to an ultrafilter U on X, then the trace of U on A is the set \(\mathfrak{P}(A)\) .
\item
  On an infinite set, show that an elementary filter associated with a sequence all of whose terms are distinct is not an ultrafilter.
\item
  Let \(f\) be a mapping of a set \(\mathbf{X}\) into a set \(\mathbf{X}'\) . Then \(f\) is injective if and only if, for every filter base \(\mathfrak{B}\) on \(\mathbf{X}\) , \((f\overline{f}^{1}(\mathfrak{B}))\) is a filter base equivalent to \(\mathfrak{B}\) .
\item
  Show that if \(f\) is a mapping of a set \(\mathbf{X}\) onto a set \(\mathbf{X}'\) , then the image \(f(\mathfrak{F})\) of every filter \(\mathfrak{F}\) on \(\mathbf{X}\) is a filter on \(\mathbf{X}'\) .
\end{enumerate}

¶ 15) Let \(n \to f(n)\) be a mapping of \(\mathbf{N}\) onto \(\mathbf{N}\) such that \(\overline{f}(m)\) is finite for each \(m \in \mathbf{N}\) . If \((x_n)\) is any sequence in a set \(\mathbf{X}\) , put \(y_n = x_{f(n)}\) . Show that the elementary filters associated with the sequences \((x_n)\) and \((y_n)\) are the same.

Deduce that if \((a_{n})\) and \((b_{n})\) are two sequences of elements of a set X such that the filter associated with \((b_{n})\) is finer than the filter associated with \((a_{n})\) , then the filter associated with \((b_{n})\) is the same as the filter associated with some subsequence of \((a_{n})\) .

¶ 16) Let \(\Phi\) be a countable linearly ordered set of elementary filters on a set X. Show that there is an elementary filter which is finer than all the filters of \(\Phi\) (show that the union of the filters of \(\Phi\) has a countable base).

\begin{enumerate}
\def\labelenumi{\arabic{enumi})}
\setcounter{enumi}{16}
\tightlist
\item
  Let X be a lattice (Set Theory, Chapter III, § 1, no. 13). A non-empty subset F of X is called a prefilter if it satisfies the following conditions: (i) whenever \(x \in \mathbf{F}\) and \(y \geqslant x\) , then \(y \in \mathbf{F}\) ; (ii) whenever \(x \in \mathbf{F}\) and \(y \in F\) , then \(\inf (x, y) \in F\) ; (iii) \(F \neq X\) . Thus the prefilters on the set of subsets of a set Y, ordered by inclusion, are precisely the filters on Y. A prefilter F is said to be prime if the relation \(\sup(x, y) \in F\) implies that either \(x \in F\) or \(y \in F\) . A prefilter with respect to the ordering opposite to that of X is called a coprefilter on X.
\end{enumerate}

\begin{enumerate}
\def\labelenumi{\alph{enumi})}
\item
  The set of upper bounds of a subset A of X which does not consist only of the least element of X is a prefilter. Give examples of prefilters which are not sets of upper bounds.
\item
  Show that if F is a prime prefilter, then CF is a prime coprefilter.
\item
  Show that if X has a least element 0, then every prefilter is contained in a maximal prefilter (use Zorn's lemma, noting that 0 belongs to no prefilter).
\item
  Find all prefilters on a linearly ordered set X, and show that they are all prime. Hence give examples of prime prefilters which are not maximal.
\item
  Let X be an ordered set of 5 elements, denoted by 0, 1, a, b, c, the order relations being 0 \textless{} a \textless{} 1, 0 \textless{} b \textless{} 1, 0 \textless{} c \textless{} 1. Show that X is a lattice and that the prefilters other than \(\{1\}\) are maximal but not prime.
\item
  Let X be a lattice having a least element 0. Show that if \(F \subset X\) is a prefilter and a is an element of X which does not belong to F, then there is a prefilter containing F and a if and only if \(\inf (a, x) \neq 0\) for all \(x \in F\) ; the smallest prefilter containing F and a is then the set of all \(y \in X\) such that \(y \geqslant \inf (a, x)\) for at least one element x of F.
\end{enumerate}

¶ 18) Let X be a distributive lattice (Set Theory, Chapter III, § 1, Exercise 16) which has a least element 0.

\begin{enumerate}
\def\labelenumi{\alph{enumi})}
\item
  If \(a\) is an irreducible element of \(\mathbf{X}\) (Set Theory, Chapter III, § 4, Exercise 7), show that the set of all upper bounds of \(a\) is a prime prefilter. Give an example of a prime prefilter which is not a set of upper bounds of an irreducible element of \(\mathbf{X}\) {[}cf.~Exercise 17 d){]}.
\item
  Show that every maximal prefilter on X is prime {[}use Exercise 17f){]}. Give examples of prime non-maximal prefilters on a distributive lattice {[}Exercise 17d){]}.
\item
  Show that every prefilter F is the intersection of the prime prefilters which contain F {[}if \(a \notin F\) , show that a maximal element U in the set of prefilters which contain F but not \(a\) is prime, by using Exercise 17f){]}.
\item
  Let \(\Omega\) be the set of prime prefilters on X. To each \(x \in \mathbf{X}\) let correspond the subset \(\mathbf{A}_x\) of \(\Omega\) consisting of prime prefilters F such that \(x \in \mathbf{F}\) . Show that the mapping \(x \to \mathbf{A}_x\) of X into \(\mathfrak{P}(\Omega)\) is injective and order-preserving, that \(\mathbf{A}_{\mathrm{inf}(x,y)} = \mathbf{A}_x \cap \mathbf{A}_y\) and that
\end{enumerate}

\[
\mathbf {A} _ {\sup (x, y)} = \mathbf {A} _ {x} \cup \mathbf {A} _ {y}.
\]

\begin{enumerate}
\def\labelenumi{\arabic{enumi})}
\setcounter{enumi}{18}
\item
  Let X be a lattice which has a least element. Show that if every prefilter on X is the intersection of the prime prefilters which contain it, then X is distributive (if the mapping \(x \rightarrow A_{x}\) is defined as in Exercise 18 d), this mapping still has the properties stated there).
\item
  A lattice X is said to be a Boolean algebra if it is distributive, has a least element 0, and a greatest element 1 and if, for each \(x \in \mathbf{X}\) , there is an element \(x' \in \mathbf{X}\) such that \(\inf (x, x') = 0\) and \(\sup (x, x') = 1\) ; such an element is said to be a complement of \(x\) in \(\mathbf{X}\) .
\end{enumerate}

\begin{enumerate}
\def\labelenumi{\alph{enumi})}
\item
  Let \(x \to \mathbf{A}_x\) be the injection of a Boolean algebra \(\mathbf{X}\) into the set of subsets of the set \(\Omega\) of prime prefilters on \(\mathbf{X}\) , defined in Exercise 18 d). Show that if \(x'\) is a complement of \(x\) in \(\mathbf{X}\) , then \(\mathbf{A}_{x'} = \complement \mathbf{A}_x\) in \(\Omega\) ; deduce that an element \(x\) has only one complement \(x'\) , that the complement of \(x'\) is \(x\) , and that the complement of \(\inf (x, y)\) is \(\sup (x', y')\) ; give direct proofs of these statements. The ideals of the Boolean ring \(\mathbf{A}\) corresponding to \(\mathbf{X}\) , other than \(\mathbf{A}\) , are the coprefilters on \(\mathbf{X}\) .
\item
  Show that, in a Boolean algebra \(\mathbf{X}\) , every prime prefilter is maximal (remark that, for each \(x \in \mathbf{X}\) , every prime prefilter must contain either \(x\) or its complement).
\item
  Conversely, let X be a distributive lattice which has a least element 0 and a greatest element 1, and is such that every prime prefilter on X is maximal. Show that X is a Boolean algebra. (For an element \(x \in X\) , consider the coprefilter C (Exercise 17) consisting of those \(y \in X\) such that \(\inf(x, y) = 0\) ; if \(x\) has no complement, there is a maximal coprefilter M containing \(x\) and C; remark that the complement U of M in X is a prime prefilter {[}Exercise 17 b) and 18 b){]} and use Exercise 17 f)). d) Show that if X is a topological space then the set B of subsets of X which are both open and closed is a Boolean algebra if ordered by inclusion. If we take X to be the topological space associated with the Fréchet filter (no. 5), show that the Boolean algebra B thus obtained is countably infinite and therefore cannot be isomorphic to the ordered set P(A) of all subsets of a set A.
\end{enumerate}
